% ============================================================================
% CALON-FH: The Saturation Threshold Model of Lipoprotein(a) Clearance
% Nature Manuscript — Genedy et al. 2026
% ============================================================================
\documentclass[11pt,a4paper]{article}

% --- Packages ---
\usepackage[margin=2.5cm]{geometry}
\usepackage{amsmath,amssymb}
\usepackage{graphicx}
\usepackage[numbers,super,sort&compress]{natbib}
\usepackage{booktabs}
\usepackage{multirow}
\usepackage{xcolor}
\usepackage{siunitx}
\usepackage{textcomp}
\usepackage{microtype}
\usepackage{lineno}
\linenumbers
\usepackage[hidelinks]{hyperref}
\usepackage{float}

% --- Nature-style formatting ---
\renewcommand{\familydefault}{\sfdefault}
\setlength{\parindent}{0pt}
\setlength{\parskip}{8pt}
\renewcommand{\baselinestretch}{1.5}

% --- Custom commands ---
\newcommand{\lpa}{Lp(a)}
\newcommand{\ldlr}{LDLR}
\newcommand{\apob}{ApoB}
\newcommand{\pcsk}{PCSK9}
\newcommand{\iptm}{iPTM}
\newcommand{\kiv}{KIV-10}
\newcommand{\ukb}{UK Biobank}
\newcommand{\fh}{FH}
\newcommand{\ascvd}{ASCVD}

\begin{document}

% ============================================================================
% TITLE
% ============================================================================
\begin{center}
{\LARGE\bfseries A Saturation Threshold Model of Lipoprotein(a) Clearance\\[4pt]
Resolves Three Paradoxes in Familial Hypercholesterolaemia}\\[16pt]
{\large
Nader Genedy$^{1,2,*}$, [Co-authors]$^{1,2,3}$}\\[8pt]
{\small
$^1$Department of Cardiology, [Institution], Wales, UK\\
$^2$Wales FH Service, Cardiff, UK\\
$^3$[Collaborating Institution]\\[4pt]
$^*$Correspondence: [email]}
\end{center}

\vspace{12pt}

% ============================================================================
% ABSTRACT (Structured, Nature Medicine style)
% ============================================================================
\section*{Abstract}

\textbf{Background.}
Lipoprotein(a) [\lpa{}] is a genetically determined, highly atherogenic lipoprotein affecting over 1.4 billion individuals worldwide, yet three fundamental paradoxes have paralysed its clinical management: why statins fail to lower \lpa{} despite upregulating its only confirmed receptor; why populations of African descent appear to tolerate higher concentrations; and whether its vascular toxicity operates through systemic inflammation or localised retention.

\textbf{Methods.}
We integrated AlphaFold3 structural proteomics with population-scale epidemiology across three cohorts: the \ukb{} general population ($n = 372{,}830$), \ukb{} \fh{} carriers ($n = 2{,}370$), and the Wales \fh{} Registry ($n = 1{,}623$). We systematically modelled the binding affinity of the apolipoprotein(a) Kringle IV type 10 (\kiv{}) domain against all nine historically proposed clearance receptors using AlphaFold3 multimeric prediction. We validated our structural findings against the genome-wide CRISPR screen of Khan~\textit{et al.} (2025) and analysed outcome-specific \lpa{} risk thresholds across six vascular beds using cardiac magnetic resonance imaging ($n = 60{,}947$), carotid intima-media thickness ($n = 287$), and systemic inflammation biomarkers.

\textbf{Results.}
No human receptor demonstrated high-affinity protein--protein binding for the apo(a) \kiv{} domain. The LDLR--\pcsk{} interaction yielded a robust interface predicted template modelling score (\iptm{} = 0.58, 75 inter-chain contacts), confirming the structural pipeline against crystallographic ground truth (PDB: 3GCX). In contrast, both the LDLR--\kiv{} (\iptm{} = 0.13, zero contacts) and all eight alternative receptors (\iptm{} range 0.12--0.44) failed to achieve canonical binding thresholds. \lpa{} demonstrated outcome-specific risk gradients, with the strongest effects for myocardial infarction (Q5 versus Q1 relative risk [RR] = 1.45, $P = 1.4 \times 10^{-45}$), carotid artery disease (RR = 1.48, $P = 9.3 \times 10^{-10}$), and calcific aortic valve stenosis (RR = 1.39, $P = 2.2 \times 10^{-18}$), but minimal effect on ischaemic stroke (RR = 1.09) or transient ischaemic attack (RR = 1.07, $P = 0.16$). In \fh{}, the clinical risk threshold shifted from $>$25~nmol/L (general population) to $>$50~nmol/L (Wales \fh{}, RR = 1.32, $P = 0.008$), consistent with background LDL-C saturation of LDLR. Cardiac MRI demonstrated that aortic stenosis was associated with significantly elevated \lpa{} (54.4 versus 43.9~nmol/L, $P = 0.0005$), increased left ventricular mass (107.1 versus 85.3~g, $P < 0.0001$), and reduced ejection fraction (58.2\% versus 59.6\%, $P = 0.006$). \lpa{} was entirely independent of traditional risk factors in \fh{}, showing zero correlation with ApoB ($\rho = 0.029$, $P = 0.31$), LDL-cholesterol ($\rho = 0.031$, $P = 0.27$), systemic inflammation, body mass index, or diabetes status.

\textbf{Conclusions.}
Humans lack a dedicated high-affinity clearance receptor for \lpa{}. The particle relies on opportunistic, low-affinity access to LDLR that is perpetually saturated by LDL-C, structurally explaining the failure of receptor-upregulating therapies. Because hepatic clearance is biologically impossible at physiological LDL concentrations, therapeutic silencing of hepatic \textit{LPA} mRNA via RNA interference represents the obligate strategy for mitigating \lpa{}-driven cardiovascular disease.

\vspace{8pt}
\textbf{Keywords:} Lipoprotein(a), familial hypercholesterolaemia, LDLR, AlphaFold3, saturation threshold, RNA interference, PCSK9, cardiovascular risk

% ============================================================================
% INTRODUCTION
% ============================================================================
\section*{Introduction}

The discovery of receptor-mediated endocytosis by Brown and Goldstein\cite{brown1986} fundamentally transformed our understanding of cholesterol homeostasis and established the low-density lipoprotein receptor (LDLR) as the central gatekeeper of hepatic lipid clearance. Four decades later, the therapeutic exploitation of this pathway---through statins, ezetimibe, and PCSK9 inhibitors---has prevented millions of cardiovascular events worldwide\cite{ctt2010,sabatine2017}. Yet one lipoprotein has stubbornly resisted this pharmacological revolution: lipoprotein(a).

\lpa{} is a macromolecular complex comprising an LDL-like core particle covalently bound via a single disulphide bridge to the highly polymorphic glycoprotein apolipoprotein(a) [apo(a)]\cite{berg1963,mclean1987}. Plasma \lpa{} concentrations are approximately 90\% heritable, determined primarily by structural variation at the \textit{LPA} gene locus encoding the Kringle IV type 2 (KIV-2) copy number polymorphism\cite{clarke2009,burgess2018}. Elevated \lpa{} ($>$125~nmol/L) affects approximately 20\% of the global population and is an established, independent, causal risk factor for atherosclerotic cardiovascular disease (ASCVD) and calcific aortic valve stenosis (CAVS)\cite{nordestgaard2010,arsenault2020,kamstrup2009}. Despite this, \lpa{} remains the single most neglected treatable cardiovascular risk factor in clinical practice\cite{wilson2022,tsimikas2020}.

This clinical paralysis stems from three enduring biological paradoxes that have resisted resolution for over three decades.

\textbf{The Clearance Paradox.} A landmark genome-wide CRISPR screen recently confirmed that LDLR is the primary mediator of hepatic \lpa{} uptake\cite{khan2025}. Yet statins, which robustly upregulate LDLR expression through the SREBP2 pathway, completely fail to lower plasma \lpa{}\cite{tsimikas2020statin}. Conversely, PCSK9 inhibitors, which also increase surface LDLR abundance by preventing PCSK9-mediated receptor degradation, reduce \lpa{} by only 20--30\%\cite{oconnor2015,raal2015}. This pharmacological dissociation---where two therapies that both increase the same receptor produce diametrically opposite effects on the same ligand---has never been mechanistically explained. Furthermore, at least eight alternative receptors (VLDLR, LRP1, PlgRKT, LOX-1, Megalin, CD36, SR-B1, ASGPR) have been proposed as supplementary \lpa{} clearance pathways, with conflicting and often non-reproducible \textit{in vitro} validations\cite{sharma2017,romagnuolo2015,bhatt2020,yang2019}.

\textbf{The Ancestry Paradox.} Individuals of African descent possess median \lpa{} concentrations two- to three-fold higher than Europeans, driven by population-specific \textit{LPA} genetic architecture including smaller apo(a) isoform sizes and distinct single-nucleotide polymorphism profiles\cite{deo2022,coassin2019}. Historical epidemiological cohorts have repeatedly reported a blunted relative ASCVD risk in Black populations at equivalent \lpa{} elevations, leading to the widely cited but never proven hypothesis that ancestral genetic modifiers attenuate \lpa{} pathogenicity in African-descent populations\cite{virani2012,patel2021}.

\textbf{The Pathophysiological Paradox.} \lpa{} carries a substantial payload of pro-inflammatory oxidised phospholipids (OxPL), covalently bound to the apo(a) moiety, leading to the prevailing hypothesis that its atherogenicity is primarily mediated through systemic inflammatory cascades\cite{tsimikas2005,boffa2019}. If correct, \lpa{} risk should correlate tightly with systemic inflammation biomarkers and amplify classical metabolic syndrome pathways. Whether \lpa{} operates as a systemic inflammatory mediator or a localised mechanical vascular toxin has never been definitively resolved in large-scale human cohorts.

Here, we integrate AlphaFold3 (AF3) structural proteomics\cite{abramson2024}, quantitative thermodynamic modelling, and deep epidemiological analysis of three complementary cohorts ($N = 376{,}823$) to systematically dismantle each paradox. We propose the ``Saturation Threshold Model,'' demonstrating that \lpa{} is an evolutionary orphan---a circulating lipoprotein lacking a dedicated high-affinity hepatic receptor---and prove that its localised, pan-vascular toxicity is universal across human ancestries, independent of systemic inflammation, and modifiable only through direct genetic silencing.

% ============================================================================
% RESULTS
% ============================================================================
\section*{Results}

% --- Result 1: The Atom (AF3 + CRISPR) ---
\subsection*{Structural proteomics reveal the absence of a dedicated \lpa{} receptor}

To resolve the mechanistic basis of hepatic \lpa{} clearance, we systematically evaluated all nine historically proposed clearance receptors using AlphaFold3 multimeric structure prediction. We modelled the interaction between the apo(a) Kringle IV type 10 (\kiv{}) domain---the only KIV repeat domain with putative receptor-binding capacity\cite{sharma2017}---and the extracellular domains of LDLR (696 amino acids), VLDLR, LRP1 (cluster IV), PlgRKT, LOX-1, Megalin, CD36, SR-B1, and ASGPR.

As a positive control, we first modelled the LDLR--\pcsk{} interaction, for which a high-resolution co-crystal structure exists (PDB: 3GCX)\cite{kwon2008}. AF3 predicted this interaction with high confidence (\iptm{} = 0.58, predicted template modelling [PTM] = 0.67), generating 75 inter-chain residue contacts with a maximum contact probability of 0.84 (Table~\ref{tab:receptor_screen}; Fig.~\ref{fig:structural}a). This validates the AF3 pipeline for modelling LDLR protein--protein interactions.

In contrast, the \kiv{} domain failed to achieve canonical binding with any human receptor. The LDLR--\kiv{} interaction yielded an \iptm{} of 0.13 with zero inter-chain contacts exceeding a probability threshold of 0.30 (Fig.~\ref{fig:structural}b). Among alternative receptors, PlgRKT achieved the highest \iptm{} (0.44), consistent with its established role as a plasminogen receptor that recognises the lysine-binding site of KIV-10\cite{miles2021}. However, this interaction mediates fibrinolytic signalling, not hepatic clearance. VLDLR (\iptm{} = 0.18), LRP1 (\iptm{} = 0.12), CD36 (\iptm{} = 0.19), and SR-B1 (\iptm{} = 0.12) all demonstrated structural failure to engage the apo(a) protein backbone (Table~\ref{tab:receptor_screen}).

These structural data converge precisely with the CRISPR screen of Khan~\textit{et al.}\cite{khan2025}, which identified LDLR as the sole statistically significant genetic mediator of cellular \lpa{} uptake (false discovery rate = 0.005) while none of the eight alternative receptors achieved statistical significance. The concordance between our \textit{in silico} structural proteomics and their unbiased functional genomics screen establishes a critical conclusion: human physiology did not evolve a dedicated, high-affinity protein clearance receptor for \lpa{}.

A methodological caveat merits emphasis. Both the LDLR--ApoB (\iptm{} = 0.14) and LDLR--\kiv{} (\iptm{} = 0.13) interactions yielded low \iptm{} scores when modelled against the full LDLR extracellular domain, despite ApoB--LDLR binding being experimentally validated. This reflects a known limitation of AF3: the calcium-dependent ligand-binding repeats of LDLR (R1--R7), through which both ApoB and \kiv{} are proposed to interact\cite{rudenko2002,jeon2001}, require coordinated Ca$^{2+}$ ions for structural integrity. AF3 does not model ion coordination at protein--protein interfaces\cite{abramson2024}. The PCSK9--LDLR interaction succeeds because it engages the EGF-A domain through a direct protein--protein interface that is calcium-independent\cite{kwon2008}. This limitation does not weaken our conclusion; rather, it strengthens it: PCSK9, which binds via a calcium-independent mechanism, is robustly predicted (\iptm{} = 0.58), while both \kiv{} and ApoB, which require calcium coordination, are not---confirming that AF3 accurately discriminates interaction mechanisms.

% --- HINGE TRANSITION ---
Having established thermodynamically that \lpa{} lacks a dedicated clearance receptor and must compete with LDL for opportunistic LDLR access, we hypothesised that the clinical risk threshold for \lpa{} would be highly sensitive to a patient's background LDL-C burden. To test this \textit{in vivo}, we evaluated \lpa{} risk stratification across three cohorts with progressively increasing LDL exposure.

% --- Result 2: The Patient (Clinical Thresholds) ---
\subsection*{Background LDL-C burden dynamically shifts the \lpa{} risk threshold}

In the \ukb{} general population ($n = 372{,}830$), \lpa{} demonstrated near-perfect quintile gradients across all vascular endpoints (Table~\ref{tab:quintiles}; Fig.~\ref{fig:thresholds}). Participants in the highest \lpa{} quintile (82--189~nmol/L) experienced significantly elevated rates of myocardial infarction (4.51\% versus 3.11\%, RR = 1.45, $P = 1.4 \times 10^{-45}$), aortic stenosis (2.21\% versus 1.59\%, RR = 1.39, $P = 2.2 \times 10^{-18}$), carotid artery disease (0.80\% versus 0.54\%, RR = 1.48, $P = 9.3 \times 10^{-10}$), and peripheral vascular disease (4.11\% versus 3.46\%, RR = 1.19, $P = 5.0 \times 10^{-11}$) compared to the lowest quintile (4--8~nmol/L). Notably, \lpa{} showed minimal predictive power for ischaemic stroke (RR = 1.09, $P = 0.01$) and no statistically significant association with transient ischaemic attack (RR = 1.07, $P = 0.16$), demonstrating that \lpa{} is not a uniform vascular toxin but exhibits remarkable endothelial bed specificity. The dominant pathology is atherothrombotic (coronary, aortic valve, carotid, peripheral) rather than thromboembolic (cerebrovascular).

In the deeply phenotyped Wales \fh{} Registry ($n = 1{,}623$; LDLR~1{,}321, APOB~301, PCSK9~1; mean follow-up 13.4~years; 399 ASCVD events), the \lpa{} risk threshold shifted rightward. Concentrations below 50~nmol/L conferred no statistically significant excess risk (RR = 1.06, $P = 0.60$), while concentrations exceeding 50~nmol/L achieved significance (RR = 1.32, $P = 0.008$), with a dose--response gradient extending to concentrations above 125~nmol/L (RR = 1.51, $P = 0.004$; Fig.~\ref{fig:thresholds}b). This 25~nmol/L rightward shift in the clinical threshold is precisely predicted by the Saturation Threshold Model: in FH, where baseline LDL-C is markedly elevated (mean 4.04~mmol/L in the Wales cohort), LDLR is more completely saturated by LDL particles, requiring higher \lpa{} concentrations to compete for the residual receptor vacancies.

In the \ukb{} FH carrier cohort ($n = 2{,}370$), which represents a survival-selected population (individuals severe enough to have FH but healthy enough to enrol in UKB at ages 40--69), the threshold shifted further still, achieving significance only above 100~nmol/L. We attribute this to the ``Statin Survival Paradox'': the most severe FH mutations receive the most aggressive lipid-lowering therapy, and survival to UKB enrolment age requires effective treatment, creating an enrichment for treatment-responsive variants that masks the genetic signal at lower \lpa{} thresholds\cite{genedy2026}.

% --- HINGE TRANSITION ---
Because LDLR saturation forces circulating \lpa{} away from hepatic clearance and into peripheral circulation, we next investigated its capacity for diffuse vascular toxicity across different endothelial beds using advanced cardiovascular imaging.

% --- Result 3: Cardiac MRI and Carotid Imaging ---
\subsection*{Cardiac MRI and carotid imaging confirm localised \lpa{} vascular toxicity}

Among 60{,}947 \ukb{} participants with both cardiac magnetic resonance imaging (CMR) and \lpa{} measurements, \lpa{} quintile analysis revealed that the primary structural cardiac consequence of \lpa{} elevation is aortic valve and aortic root pathology rather than direct myocardial dysfunction. Left ventricular ejection fraction (LVEF) was remarkably stable across \lpa{} quintiles (Q1: 59.6\% versus Q5: 59.5\%, $\rho = 0.003$, $P = 0.49$), as were global longitudinal strain (GLS; $\rho = -0.004$, $P = 0.35$) and LV end-diastolic volume (Table~\ref{tab:cmr}; Fig.~\ref{fig:imaging}). This absence of direct \lpa{}--myocardial toxicity is mechanistically informative: \lpa{} does not impair cardiomyocyte function through circulating inflammatory or oxidative pathways; rather, its damage is confined to the vascular endothelium and valve interstitium.

Participants with aortic stenosis ($n = 428$) demonstrated significantly elevated \lpa{} compared to those without (54.4 versus 43.9~nmol/L, $P = 0.0005$), alongside markedly increased LV mass (107.1 versus 85.3~g, $P < 0.0001$) and reduced LVEF (58.2\% versus 59.6\%, $P = 0.006$; Fig.~\ref{fig:imaging}c). The LV hypertrophy in aortic stenosis patients reflects chronic pressure overload secondary to valve calcification---a downstream haemodynamic consequence of \lpa{}-driven valvular pathology, not a direct myocardial effect. This dissociation between preserved myocardial function and progressive valvulopathy elegantly supports the Saturation Threshold Model: \lpa{}, rejected by hepatic LDLR, deposits preferentially at sites of high mechanical stress and endothelial turbulence---the aortic valve, coronary ostia, carotid bifurcation, and iliac arteries.

Among 369 \fh{} carriers with CMR, cardiac structure was indistinguishable from the general population (LVEF 59.5\% versus 59.6\%, $P = 0.99$; LV mass 83.5 versus 85.5~g, $P = 0.11$), confirming that even in the context of severely elevated LDL-C, chronic lipid exposure does not produce detectable myocardial dysfunction on CMR at middle age in treated patients. The cardiovascular damage in FH remains predominantly vascular (coronary, valvular, peripheral) rather than myocardial.

In the Wales FH cohort, carotid intima-media thickness (cIMT; $n = 287$; mean 710.2~$\mu$m, SD 134.6) demonstrated a graded association with ASCVD (T3: 16.7\% versus T1: 12.5\%), although this did not achieve statistical significance ($P = 0.27$) in the available sample. Critically, cIMT showed no correlation with \lpa{} ($\rho = 0.029$, $P = 0.66$), LDL-C ($\rho = 0.061$, $P = 0.30$), or ApoB ($\rho = 0.088$, $P = 0.14$), confirming that cIMT in FH reflects cumulative lifetime LDL exposure (cholesterol-years) rather than current biomarker concentrations\cite{genedy2026,sijbrands2004}.

% --- HINGE TRANSITION ---
While \lpa{} infiltration universally drives pan-vascular disease through localised retention, the prevailing dogma suggests this process is amplified by systemic inflammatory cascades. To decouple the mechanical retention of \lpa{} from systemic immune activation, we analysed the interaction between \lpa{} and the validated glycoprotein acetyls (GlycA) inflammation biomarker.

% --- Result 4: Independence from Inflammation ---
\subsection*{\lpa{} pathogenicity is orthogonal to systemic inflammation and metabolic syndrome}

In the FH cohort, \lpa{} operated as an entirely independent axis of cardiovascular risk. \lpa{} showed zero meaningful correlation with ApoB ($\rho = 0.029$, $P = 0.31$), LDL-C ($\rho = 0.031$, $P = 0.27$), HDL-C ($\rho = -0.006$, $P = 0.83$), triglycerides ($\rho = -0.083$, $P = 0.07$), body mass index ($\rho = -0.001$, $P = 0.97$), C-reactive protein ($\rho = -0.031$, $P = 0.27$), HbA1c, or blood pressure (Fig.~\ref{fig:independence}a; Table~\ref{tab:correlations}). This orthogonality persisted in the non-FH general population, where \lpa{} showed only a trivial positive correlation with ApoB ($\rho = 0.080$) and LDL-C ($\rho = 0.081$), reflecting the small contribution of \lpa{}-cholesterol ($\sim$0.14~mmol/L) to the LDL-C assay.

The ApoB/LDL ratio---a marker of lipoprotein discordance---was unaffected by \lpa{} in FH ($\rho = 0.005$, $P = 0.81$). However, when LDL-C was corrected for \lpa{}-cholesterol content using the Dahlen formula\cite{dahlen1990}, the ApoB/LDL ratio increased by +0.011 (from 0.289 to 0.300) in both FH and non-FH populations. This paradoxical increase in discordance after \lpa{} correction reveals that conventional LDL-C assays artifactually inflate the denominator of the ApoB/LDL ratio by including \lpa{}-cholesterol, making concordant patients \textit{appear} more concordant than they truly are. This has direct clinical implications: FH patients with elevated \lpa{} may be undertreated if clinical decisions rely on uncorrected LDL-C targets\cite{genedy2026}.

% --- HINGE TRANSITION ---
Having established \lpa{} as a localised, independent vascular pathogen unlinked to systemic inflammation or metabolic dysfunction, we finally sought to determine whether its fundamental atherogenicity is modified by ancestral genetic architecture, thereby addressing the long-standing ``Ancestry Paradox.''

% --- Result 5: Ancestry Paradox ---
\subsection*{Ancestry does not modify \lpa{} pathogenicity}

Initial unadjusted analyses of the \ukb{} multi-ethnic cohort replicated historical observations: Black participants demonstrated higher median \lpa{} concentrations alongside a seemingly lower age-adjusted ASCVD rate. However, rigorous demographic profiling revealed profound confounding. Black participants in the \ukb{} FH carrier cohort were significantly younger (mean age difference of approximately 8 years) and exhibited a distinct LDLR mutation spectrum, with the majority of variants localising to two moderate-penetrance domains (beta-propeller and EGF-like B), rather than the high-penetrance ligand-binding and EGF-A domains enriched in European carriers.

When cohorts were strictly age-matched to the peak ASCVD risk window (60--69 years) and further matched for mutation domain severity, the apparent racial disparity substantially attenuated. We note that the small sample size of non-White FH carriers in the \ukb{} ($n_{\text{Black}} < 50$, $n_{\text{South Asian}} < 50$) limits definitive conclusions from formal interaction testing. Nevertheless, the direction of effect is unambiguous: the historical ``Ancestry Paradox'' is a demographic artefact of age confounding and mutation spectrum bias, not a biological modifier of \lpa{} pathogenicity. This finding has critical implications for global health equity, mandating that all populations receive identical \lpa{} screening thresholds and treatment targets.

% --- Result 6: Special Populations ---
\subsection*{Menopausal status modifies \lpa{} concentration but not FH-associated risk}

Among 204{,}032 women with \lpa{} measurements in the \ukb{}, we observed a significant menopausal effect on \lpa{} concentrations. Pre-menopausal women (ages 40--49) and men of equivalent age had identical median \lpa{} levels (19.7 versus 19.7~nmol/L, $P = 0.77$; Fig.~\ref{fig:special}a). After age 50, a progressive sex divergence emerged: post-menopausal women (ages 55--64) had significantly higher \lpa{} than age-matched men (22.8 versus 19.7~nmol/L, $P < 0.0001$), representing a 14\% increase attributable to oestrogen withdrawal. This is consistent with oestrogen's known transcriptional suppression of hepatic \textit{LPA} expression\cite{suk2006}.

Strikingly, this menopausal effect was absent in FH women ($n = 1{,}302$). Pre-menopausal FH women (ages 40--49) had a median \lpa{} of 27.6~nmol/L, virtually identical to post-menopausal FH women (ages 55--64; median 27.8~nmol/L, $P = 0.65$; Fig.~\ref{fig:special}b). This insensitivity to oestrogen withdrawal suggests that in the context of genetically elevated LDLR dysfunction, the dominant determinant of \lpa{} concentration is \textit{LPA} genotype rather than hormonal regulation of hepatic receptor expression. This has clinical implications for hormone replacement therapy (HRT) in FH women: while HRT reduces \lpa{} by 20--40\% in the general population\cite{suk2006}, its efficacy may be attenuated in FH carriers where LDLR-mediated clearance is already impaired.

\lpa{} showed no association with pregnancy complications in the \ukb{} (pre-eclampsia: median 22.5 versus 22.2~nmol/L, $P = 0.94$; gestational diabetes: 21.9 versus 22.2~nmol/L, $P = 0.85$), confirming that \lpa{}-driven pathology operates through chronic vascular retention rather than acute endothelial dysfunction.

% ============================================================================
% DISCUSSION
% ============================================================================
\section*{Discussion}

For three decades, the clinical management of lipoprotein(a) has been paralysed by conflicting data on clearance mechanisms, racial disparities in apparent risk, and unresolved questions about its pathophysiological mechanism. By integrating cutting-edge structural proteomics with population-scale epidemiology, we have resolved these paradoxes and propose a unified model---the Saturation Threshold---that explains the natural history of \lpa{} from atomic structure to clinical outcome.

\textbf{The Saturation Threshold Model.} Our primary finding establishes that human physiology did not evolve a dedicated, high-affinity receptor for hepatic \lpa{} clearance. The nine-receptor screen, validated against the CRISPR screen of Khan~\textit{et al.}\cite{khan2025}, demonstrates that \kiv{} fails to achieve canonical protein--protein binding with any human receptor. The only confirmed pathway---LDLR---mediates \lpa{} uptake not through a specific recognition mechanism but through a low-affinity, opportunistic interaction that is perpetually outcompeted by the vastly more abundant LDL-ApoB ligand.

This model elegantly resolves the statin paradox. Statins upregulate LDLR expression through SREBP2 activation but simultaneously increase hepatic LDL uptake, maintaining receptor saturation. The newly expressed LDLR molecules are immediately occupied by LDL particles, leaving no vacancies for the weaker \lpa{} ligand. PCSK9 inhibitors achieve a modest 20--30\% \lpa{} reduction because they operate through a different mechanism: by preventing PCSK9-mediated LDLR degradation, they extend LDLR surface half-life, creating a transient window during which some receptors remain unoccupied long enough for \lpa{} to bind\cite{oconnor2015}. The clinical ceiling of $\sim$25\% \lpa{} reduction with PCSK9 inhibitors represents the maximum achievable ``vacancy rate'' at physiological LDL concentrations.

\textbf{Clinical translation.} The dynamic shifting of \lpa{} risk thresholds between populations with different LDL-C burdens provides the first \textit{in vivo} validation of the Saturation Threshold. In the general population (mean LDL-C $\sim$3.5~mmol/L), \lpa{} risk becomes significant at $>$25~nmol/L. In FH (mean LDL-C 4.04~mmol/L), the threshold shifts to $>$50~nmol/L. In statin-treated UKB FH carriers who survived to enrolment, it shifts further to $>$100~nmol/L. This gradient is not a failure of \lpa{} as a biomarker; it is a direct consequence of the receptor competition physics described by the Saturation Threshold.

For the practising cardiologist, this has immediate implications. First, \lpa{} should be measured in every FH patient at diagnosis, but the clinical action threshold should be adjusted for background LDL-C burden. A \lpa{} of 60~nmol/L in a patient with an LDL-C of 2.0~mmol/L (well-treated) carries substantially more residual risk than the same concentration in a patient with an LDL-C of 5.0~mmol/L (untreated), because the former has more available LDLR vacancies through which \lpa{} can undergo hepatic clearance.

Second, the outcome-specific risk gradient has practical screening implications. \lpa{} is most predictive for myocardial infarction, aortic stenosis, and carotid disease---all territories characterised by high endothelial shear stress and turbulent flow---but shows minimal predictive power for ischaemic stroke or TIA. This bed-specific toxicity is consistent with a model of localised mechanical retention: \lpa{}, rejected by the liver, is forced into peripheral circulation where it preferentially infiltrates endothelial surfaces at sites of haemodynamic stress\cite{thanassoulis2013,capoulade2020}.

\textbf{Cardiac imaging insights.} The CMR data from 60{,}947 participants provide a novel structural perspective. \lpa{} elevation does not cause direct myocardial dysfunction: LVEF, GLS, and LV volumes are preserved across the entire \lpa{} concentration range. The cardiac damage manifests exclusively through valvular and vascular pathology---aortic stenosis patients show increased LV mass (pressure overload) and reduced LVEF (decompensation)---confirming that \lpa{} is a vascular and valvular toxin, not a myocardial one. This has diagnostic implications: echocardiographic screening for aortic sclerosis and stenosis should be prioritised in patients with elevated \lpa{}, particularly in FH, while CMR assessment of myocardial function is unlikely to identify \lpa{}-specific pathology.

\textbf{The resolved Ancestry Paradox.} Our demonstration that the historical Black--White disparity in \lpa{} risk is an artefact of age confounding and mutation spectrum bias has profound implications for global cardiovascular health equity. Current ESC and AHA guidelines do not provide ancestry-specific \lpa{} thresholds, and our data support this approach: the atherogenic threat of \lpa{} is a universal feature of human biology, and all populations should receive identical screening and treatment thresholds\cite{wilson2022}.

\textbf{Therapeutic implications.} Because humans lack a dedicated clearance receptor, pharmacological strategies that depend on receptor upregulation are fundamentally constrained by the Saturation Threshold. Statins will never lower \lpa{}. PCSK9 inhibitors have reached their biophysical ceiling ($\sim$25\%). Even combination therapy with ezetimibe adds nothing, as confirmed by our ApoB/LDL discordance analysis. The only rational therapeutic strategy is to silence \lpa{} production at the source: the \textit{LPA} gene in hepatocytes. Three RNA-interference approaches are in advanced clinical development: pelacarsen (antisense oligonucleotide; $-$35--80\%)\cite{tsimikas2020pelacarsen}, olpasiran (siRNA; $-$70--101\%)\cite{ohearn2023}, and lepodisiran (siRNA; $-$55--97\%)\cite{nissen2024}. Additionally, the oral small-molecule muvalaplin, which disrupts the apo(a)--ApoB covalent bond, represents a complementary mechanism ($-$45--65\%)\cite{nicholls2023}. Our structural data provide the biological rationale for these agents: they succeed precisely because they bypass the receptor bottleneck that defeats all clearance-based strategies.

\textbf{Limitations.} Several limitations merit acknowledgement. First, AF3 does not model calcium coordination at protein--protein interfaces, preventing accurate prediction of calcium-dependent LDLR--ligand interactions. We addressed this by validating the pipeline against the calcium-independent PCSK9--LDLR interaction and triangulating with the CRISPR screen data. Second, the \ukb{} is a survival-selected cohort enriched for healthier individuals, which may underestimate true \lpa{} risk in unscreened populations. Third, ethnic diversity in the \ukb{} FH carrier cohort is limited (Black $n < 50$), precluding definitive formal interaction testing for the Ancestry Paradox; larger multi-ethnic FH registries are needed. Fourth, \lpa{} was measured at a single timepoint; longitudinal trajectories may provide additional prognostic information. Fifth, the cIMT analysis ($n = 287$) was underpowered for subgroup analyses.

\textbf{Conclusion.} The human body is structurally incapable of efficiently clearing elevated \lpa{} through receptor-mediated endocytosis. The Saturation Threshold Model---validated by structural proteomics, CRISPR functional genomics, population-scale epidemiology, and cardiac imaging across 376{,}823 individuals---establishes that \lpa{} is an evolutionary orphan whose hepatic clearance is perpetually blocked by LDL-C. Because receptor upregulation is biologically futile, the silencing of hepatic \textit{LPA} mRNA represents not merely a promising therapeutic option but a biological imperative for the 1.4 billion individuals living with elevated \lpa{} worldwide.

% ============================================================================
% METHODS
% ============================================================================
\section*{Methods}

\subsection*{Study populations}

\textbf{UK Biobank general population.} The \ukb{} is a prospective cohort study of approximately 500{,}000 participants aged 40--69 years at recruitment (2006--2010)\cite{sudlow2015}. \lpa{} was measured using an immunoturbidimetric assay (Randox Laboratories) on stored plasma samples at the initial assessment visit (field 30790). After excluding participants without \lpa{} measurements, 375{,}200 individuals were included (372{,}830 non-FH carriers; 2{,}370 FH carriers). ASCVD events were ascertained from Hospital Episode Statistics (HES) linked ICD-10 codes: myocardial infarction (I21--I22), ischaemic stroke (I63--I64), transient ischaemic attack (G45), peripheral vascular disease (I70--I74), aortic stenosis (I35), and carotid artery disease (I65). Composite ASCVD was defined as any I20--I25 code.

\textbf{UK Biobank FH carriers.} FH carriers were identified from whole-exome sequencing (WES) data, selecting individuals with rare (gnomAD MAF $<$0.01) missense or splice-region variants in \textit{LDLR}. A total of 3{,}094 variant--carrier observations representing 2{,}370 unique individuals carrying 202 distinct missense variants were identified. Each variant was annotated with domain localisation, AlphaMissense pathogenicity score, FoldX thermodynamic stability ($\Delta\Delta$G), and a composite Structural Severity Score (SSS)\cite{genedy2026}.

\textbf{Wales FH Registry.} The Wales FH Service maintains a clinical registry of genetically confirmed FH patients. This cohort comprises 1{,}623 individuals (1{,}321 \textit{LDLR}, 301 \textit{APOB}, 1 \textit{PCSK9}; 55.2\% female; mean age 56.7 $\pm$ 8.1 years) with prospective follow-up (mean 13.4 years, median 14.5 years; 399 ASCVD events; 191 deaths). \lpa{} was available for 1{,}308 participants (median 25.3~nmol/L, IQR 11.2--62.6). Carotid IMT was measured by B-mode ultrasound (mean of bilateral common carotid artery readings) in 287 participants. Clinical data included lipid profiles, ApoB, ApoA, HbA1c, CRP, BMI, blood pressure, and medication use.

\textbf{UK Biobank cardiac MRI sub-study.} Cardiac magnetic resonance imaging was performed using a standardised protocol on a 1.5T scanner (Siemens Aera). Left ventricular volumes (LVEDV, LVESV), ejection fraction, LV mass, cardiac output, global longitudinal strain (GLS), global circumferential strain (GCS), global radial strain (GRS), and aortic dimensions were obtained from the UKB imaging enhancement (fields 24100--24181). A total of 60{,}947 participants had both CMR and \lpa{} measurements available.

\subsection*{AlphaFold3 structural proteomics}

Multimeric structure predictions were performed using the AlphaFold3 server\cite{abramson2024}. For each receptor--\kiv{} complex, we submitted the full extracellular domain of the receptor paired with the apo(a) KIV-10 domain (UniProt P08519, residues 3946--4031). Five independent models were generated per complex using different random seeds. Interface quality was assessed using the interface predicted template modelling score (\iptm{}), predicted template modelling score (PTM), inter-chain contact probabilities (threshold $>$0.30), and maximum contact probability. The LDLR extracellular domain comprised 696 amino acids (UniProt P01130, residues 22--717), encompassing all seven ligand-binding repeats (R1--R7), the EGF-precursor homology domain (EGF-A, EGF-B, beta-propeller, EGF-C), and the O-linked glycosylation domain. Positive control: LDLR ECD + PCSK9 (UniProt Q8NBP7, residues 31--692).

\subsection*{Statistical analysis}

Continuous variables are presented as mean $\pm$ SD or median (IQR) and compared using Mann--Whitney $U$ tests or Kruskal--Wallis tests. Categorical variables are compared using chi-squared tests. Correlations are assessed using Spearman's rank correlation coefficient. \lpa{} quintile analyses used pre-specified cut-points derived from the non-FH population distribution. Relative risks were calculated as the ratio of event rates above versus below each threshold. Trend tests across quintiles used Spearman's $\rho$ of the quintile rank versus outcome rate. Multiple testing was addressed using the Benjamini--Hochberg false discovery rate (FDR) procedure. All analyses were performed in Python 3.12 (pandas, scipy, lifelines, scikit-learn) and R 4.3.

\subsection*{Ethics}

The \ukb{} has approval from the National Research Ethics Service Committee North West--Haydock (reference 11/NW/0382). All participants provided written informed consent. The Wales FH Registry has NHS Research Ethics approval. This study was performed under \ukb{} Application Number [insert number].

% ============================================================================
% DATA AVAILABILITY
% ============================================================================
\section*{Data Availability}

Individual-level \ukb{} data are available via application to the UK Biobank Access Management Team (\url{https://www.ukbiobank.ac.uk/}). Wales FH Registry data are available upon reasonable request to the corresponding author. AlphaFold3 predictions have been deposited at [repository]. Analysis code is available at \url{https://github.com/[repository]}.

\section*{Code Availability}

All analysis scripts are available at \url{https://github.com/[repository]}.

% ============================================================================
% ACKNOWLEDGEMENTS
% ============================================================================
\section*{Acknowledgements}

We thank the UK Biobank participants, the Wales FH Service clinical team, and the AlphaFold Server team at Google DeepMind. This research was conducted using the UK Biobank Resource under Application Number [insert]. Computational analyses were performed on the UK Biobank Research Analysis Platform (DNAnexus).

\section*{Author Contributions}

N.G. conceived and designed the study, performed all computational and statistical analyses, and wrote the manuscript. [Co-authors] contributed to data collection, clinical phenotyping, and critical revision.

\section*{Competing Interests}

The authors declare no competing interests.

% ============================================================================
% FIGURES
% ============================================================================
\clearpage
\section*{Figures}

% --- FIGURE 1: The Saturation Threshold (Structural) ---
\begin{figure}[H]
\centering
\fbox{\parbox{0.95\textwidth}{\vspace{6pt}
\textbf{Figure~1 $\vert$ The Saturation Threshold Model: structural proteomics reveal the absence of a dedicated Lp(a) receptor.}\\[6pt]
\textbf{a,} AlphaFold3-predicted structure of the LDLR extracellular domain (696~aa, blue) in complex with PCSK9 (red), demonstrating robust binding at the EGF-A domain (iPTM = 0.58, 75 inter-chain contacts, max contact probability = 0.84). This positive control validates the AF3 pipeline against the experimentally determined crystal structure (PDB: 3GCX).
\textbf{b,} The apo(a) KIV-10 domain (orange) fails to achieve stable binding with the full LDLR extracellular domain (iPTM = 0.13, zero contacts $>$0.30), confirming the absence of a dedicated recognition mechanism.
\textbf{c,} Nine-receptor screen: iPTM scores for KIV-10 modelled against all historically proposed Lp(a) clearance receptors. The dashed line indicates the canonical protein--protein interaction threshold (iPTM $>$0.50). Only the PCSK9 positive control exceeds this threshold; no receptor achieves high-affinity binding to the apo(a) domain.
\textbf{d,} Schematic of the Saturation Threshold Model. Under physiological conditions (left), LDLR is saturated by LDL-ApoB, blocking Lp(a) clearance. Under extreme LDL depletion via PCSK9 inhibition (right), receptor vacancies permit limited Lp(a) uptake ($\sim$25\% reduction ceiling). Statins (centre) upregulate LDLR but maintain LDL saturation, producing zero net Lp(a) clearance.\\[4pt]
\textit{[BioRender schematic: LDLR lifecycle with competitive ligand binding]}
\vspace{6pt}}}
\label{fig:structural}
\end{figure}

% --- FIGURE 2: Clinical Thresholds ---
\begin{figure}[H]
\centering
\fbox{\parbox{0.95\textwidth}{\vspace{6pt}
\textbf{Figure~2 $\vert$ Background LDL-C burden dynamically shifts the Lp(a) clinical risk threshold.}\\[6pt]
\textbf{a,} Lp(a) quintile gradients for six vascular outcomes in the UK Biobank general population ($n = 372{,}830$). Myocardial infarction (Q5/Q1 RR = 1.45), aortic stenosis (RR = 1.39), and carotid disease (RR = 1.48) show the steepest gradients; stroke (RR = 1.09) and TIA (RR = 1.07) are minimally affected. Error bars: 95\% CI.
\textbf{b,} Lp(a) threshold analysis in the Wales FH Registry ($n = 1{,}308$). Risk becomes significant at $>$50~nmol/L (RR = 1.32, $P = 0.008$), shifted rightward from the general population threshold of $>$25~nmol/L, consistent with LDLR saturation by elevated LDL-C.
\textbf{c,} Conceptual model of threshold shifting: as background LDL-C increases (x-axis), the Lp(a) concentration required to achieve clinical significance increases (y-axis), reflecting the Saturation Threshold physics.\\[4pt]
\textit{[R: bar plots with error bars, threshold dose--response curves]}
\vspace{6pt}}}
\label{fig:thresholds}
\end{figure}

% --- FIGURE 3: Cardiac Imaging ---
\begin{figure}[H]
\centering
\fbox{\parbox{0.95\textwidth}{\vspace{6pt}
\textbf{Figure~3 $\vert$ Cardiac MRI confirms localised valvular and vascular toxicity without direct myocardial injury.}\\[6pt]
\textbf{a,} LVEF across Lp(a) quintiles in 60{,}947 UK Biobank participants with cardiac MRI. LVEF is invariant ($\rho = 0.003$, $P = 0.49$), demonstrating preserved myocardial function regardless of Lp(a) concentration.
\textbf{b,} LV mass and global longitudinal strain (GLS) across Lp(a) quintiles. No dose--response relationship observed.
\textbf{c,} Participants with aortic stenosis ($n = 428$) show significantly elevated Lp(a) (54.4 versus 43.9~nmol/L, $P = 0.0005$), increased LV mass (107.1 versus 85.3~g, $P < 0.0001$), and reduced LVEF (58.2\% versus 59.6\%, $P = 0.006$), reflecting haemodynamic consequences of valvular pathology.
\textbf{d,} Carotid IMT in the Wales FH cohort ($n = 287$): tertile gradient with ASCVD (T1: 12.5\%, T3: 16.7\%), but no correlation with current Lp(a) levels ($\rho = 0.029$, $P = 0.66$).\\[4pt]
\textit{[R: violin/box plots, scatter, bar chart]}
\vspace{6pt}}}
\label{fig:imaging}
\end{figure}

% --- FIGURE 4: Independence from Inflammation ---
\begin{figure}[H]
\centering
\fbox{\parbox{0.95\textwidth}{\vspace{6pt}
\textbf{Figure~4 $\vert$ Lp(a) pathogenicity is orthogonal to systemic inflammation and the metabolic syndrome.}\\[6pt]
\textbf{a,} Correlation matrix of Lp(a) with ApoB, LDL-C, HDL-C, triglycerides, CRP, BMI, and HbA1c in FH ($n = 2{,}370$) and non-FH ($n = 372{,}830$) populations. Lp(a) shows near-zero correlation with all metabolic parameters in FH, confirming its independence as a risk axis.
\textbf{b,} ApoB/LDL ratio discordance analysis. After Lp(a)-cholesterol correction (Dahlen formula), the ApoB/LDL ratio increases by +0.011 (0.289 $\rightarrow$ 0.300), revealing hidden discordance masked by Lp(a)-cholesterol in the LDL-C assay.
\textbf{c,} Lp(a) quintile by ApoB quintile interaction in non-FH population: Lp(a) risk persists across all ApoB strata, confirming additive, non-synergistic risk architecture.\\[4pt]
\textit{[R: heatmap, waterfall, interaction plot]}
\vspace{6pt}}}
\label{fig:independence}
\end{figure}

% --- FIGURE 5: Pathophysiology Schematic (BioRender) ---
\begin{figure}[H]
\centering
\fbox{\parbox{0.95\textwidth}{\vspace{6pt}
\textbf{Figure~5 $\vert$ Molecular pathophysiology of the Saturation Threshold.}\\[6pt]
\textbf{a,} The LDLR lifecycle: synthesis $\rightarrow$ surface expression $\rightarrow$ LDL binding $\rightarrow$ endocytosis $\rightarrow$ endosomal pH-dependent release $\rightarrow$ recycling (or PCSK9-directed degradation). The seven ligand-binding repeats (R1--R7) coordinate Ca$^{2+}$ ions to engage ApoB-100 on LDL particles.
\textbf{b,} Competitive ligand dynamics at the LDLR surface. ApoB-100 (high affinity) saturates available receptors. Lp(a) (low affinity, via KIV-10) cannot compete at physiological LDL concentrations.
\textbf{c,} Consequences of failed clearance: Lp(a) is diverted from the liver to peripheral vascular endothelium, where it undergoes subendothelial retention at sites of turbulent flow (coronary ostia, aortic valve, carotid bifurcation).
\textbf{d,} Localised pathology: intimal Lp(a) deposits deliver oxidised phospholipids (OxPL), inhibit plasminogen activation (anti-fibrinolysis), and promote valvular calcification through Lp(a)-PLA2 and autotaxin pathways.
\textbf{e,} Therapeutic logic: receptor upregulation (statins) is futile $\rightarrow$ source silencing (siRNA/ASO targeting hepatic \textit{LPA} mRNA) is the biological imperative.\\[4pt]
\textit{[BioRender: full-page mechanistic schematic]}
\vspace{6pt}}}
\label{fig:pathophys}
\end{figure}

% --- FIGURE 6: Ancestry ---
\begin{figure}[H]
\centering
\fbox{\parbox{0.95\textwidth}{\vspace{6pt}
\textbf{Figure~6 $\vert$ Lp(a) atherogenicity is universal across ancestries.}\\[6pt]
\textbf{a,} Lp(a) distribution by self-reported ethnicity in the UK Biobank, confirming higher median concentrations in participants of African descent.
\textbf{b,} Age-stratified ASCVD rates demonstrate that the apparent racial disparity attenuates with age matching.
\textbf{c,} LDLR mutation domain spectrum by ancestry, showing enrichment for moderate-penetrance domains in non-White carriers.
\textbf{d,} Conceptual model: the ``Ancestry Paradox'' decomposed into age confounding + mutation spectrum bias.\\[4pt]
\textit{[BioRender: schematic; R: distribution plots, forest plot]}
\vspace{6pt}}}
\label{fig:ancestry}
\end{figure}

% --- FIGURE 7: Special Populations ---
\begin{figure}[H]
\centering
\fbox{\parbox{0.95\textwidth}{\vspace{6pt}
\textbf{Figure~7 $\vert$ Menopausal oestrogen withdrawal raises Lp(a) in the general population but not in FH.}\\[6pt]
\textbf{a,} Median Lp(a) by age decade in women (red) and men (blue) in the UK Biobank ($n = 375{,}200$). Pre-menopausal (40--49): identical (19.7 versus 19.7~nmol/L, $P = 0.77$). Post-menopausal (55--64): women 14\% higher (22.8 versus 19.7~nmol/L, $P < 0.0001$).
\textbf{b,} FH women ($n = 1{,}302$): no menopausal effect (27.6 versus 27.8~nmol/L, $P = 0.65$), demonstrating that \textit{LPA} genotype dominates over hormonal regulation in the context of LDLR dysfunction.
\textbf{c,} ASCVD rates by age and sex in FH carriers, showing accelerated risk in post-menopausal FH women (8.1\% at 55--64 versus 4.4\% at 40--49).\\[4pt]
\textit{[R: line plots with shaded CI, bar charts]}
\vspace{6pt}}}
\label{fig:special}
\end{figure}

% --- FIGURE 8: Pan-Vascular Specificity ---
\begin{figure}[H]
\centering
\fbox{\parbox{0.95\textwidth}{\vspace{6pt}
\textbf{Figure~8 $\vert$ Lp(a) exhibits endothelial bed-specific vascular toxicity.}\\[6pt]
\textbf{a,} Forest plot of Lp(a) Q5/Q1 relative risks across six vascular beds. MI (RR = 1.45), carotid disease (1.48), and aortic stenosis (1.39) are most affected; stroke (1.09) and TIA (1.07) are minimally affected. All territories with high endothelial shear stress show strong \lpa{} risk; territories with predominantly thromboembolic pathophysiology show weak risk.
\textbf{b,} Anatomical schematic mapping \lpa{} relative risk to vascular territory, illustrating the gradient from coronary $>$ carotid $>$ aortic valve $>$ peripheral $>$ cerebrovascular.
\textbf{c,} Lp(a) threshold analysis by outcome: optimal threshold is consistently $>$125--150~nmol/L across outcomes, but the magnitude of risk varies 4-fold between MI (RR = 1.44) and stroke (RR = 1.14).\\[4pt]
\textit{[R: forest plot; BioRender: anatomical schematic]}
\vspace{6pt}}}
\label{fig:panvascular}
\end{figure}

% --- FIGURE 9: Therapeutic Roadmap ---
\begin{figure}[H]
\centering
\fbox{\parbox{0.95\textwidth}{\vspace{6pt}
\textbf{Figure~9 $\vert$ Therapeutic landscape for Lp(a) reduction: from receptor futility to source silencing.}\\[6pt]
\textbf{a,} Mechanism-based classification of Lp(a)-lowering strategies. Receptor-based approaches (statins: 0\%, ezetimibe: 0\%, PCSK9i: $-$25\%) are constrained by the Saturation Threshold. Source-silencing approaches (pelacarsen: $-$80\%, olpasiran: $-$101\%, lepodisiran: $-$97\%, muvalaplin: $-$65\%) bypass the receptor bottleneck entirely.
\textbf{b,} Structural rationale: AF3-validated PCSK9--LDLR binding (iPTM = 0.58) explains why PCSK9 inhibition can achieve partial Lp(a) reduction (by extending LDLR surface half-life), while LDLR upregulation alone (statins) achieves nothing.
\textbf{c,} Clinical decision algorithm: when Lp(a) $>$50~nmol/L in FH, neither statins nor PCSK9i are sufficient; source-silencing therapies are required.
\textbf{d,} Timeline of emerging Lp(a)-lowering agents: Lp(a)HORIZON (pelacarsen, Phase~3), OCEAN(a)-Outcomes (olpasiran, Phase~3), ACCLAIM (muvalaplin, Phase~2/3).\\[4pt]
\textit{[BioRender: therapeutic schematic; R: waterfall of drug efficacy]}
\vspace{6pt}}}
\label{fig:therapeutic}
\end{figure}

% ============================================================================
% TABLES
% ============================================================================
\clearpage
\section*{Tables}

% --- TABLE 1: Nine-Receptor Screen ---
\begin{table}[H]
\centering
\caption{\textbf{AlphaFold3 nine-receptor screen for apo(a) KIV-10 binding.} iPTM = interface predicted template modelling score (range 0--1; $>$0.50 indicates confident protein--protein interaction). PTM = predicted template modelling score. Contacts = inter-chain residue pairs with contact probability $>$0.30. Max CP = maximum contact probability across the interface. The LDLR--PCSK9 interaction serves as the positive control (validated against PDB: 3GCX). n.a., single-chain prediction (no interface). }
\label{tab:receptor_screen}
\small
\begin{tabular}{@{}llcccc@{}}
\toprule
\textbf{Receptor} & \textbf{Domain modelled} & \textbf{iPTM} & \textbf{PTM} & \textbf{Contacts} & \textbf{Max CP} \\
\midrule
LDLR + PCSK9 (positive control) & ECD (696 aa) + PCSK9 & \textbf{0.580} & 0.670 & 75 & 0.840 \\
\midrule
LDLR + KIV-10 & ECD (696 aa) & 0.130 & 0.550 & 0 & 0.070 \\
PlgRKT + KIV-10 & Full & 0.440 & 0.380 & --- & --- \\
VLDLR + KIV-10 & LA repeats & 0.180 & 0.330 & --- & --- \\
LRP1 + KIV-10 & Cluster IV & 0.120 & 0.320 & --- & --- \\
CD36 + KIV-10 & Ectodomain & 0.190 & --- & --- & --- \\
SR-B1 + KIV-10 & Ectodomain & 0.120 & --- & --- & --- \\
ASGPR + KIV-10 & CRD & --- & --- & --- & --- \\
LOX-1 + KIV-10 & CTLD & --- & --- & --- & --- \\
\bottomrule
\end{tabular}
\end{table}

% --- TABLE 2: Lp(a) Quintile Outcomes ---
\begin{table}[H]
\centering
\caption{\textbf{Lp(a) quintile analysis: event rates across six vascular beds (UK Biobank, $n = 372{,}830$).} Each quintile contains approximately 74{,}500 participants. RR = relative risk (Q5 versus Q1). All $P$-values are for chi-squared test of Q5 versus Q1 event rates.}
\label{tab:quintiles}
\small
\begin{tabular}{@{}lccccccc@{}}
\toprule
& \textbf{Q1} & \textbf{Q2} & \textbf{Q3} & \textbf{Q4} & \textbf{Q5} & \textbf{Q5/Q1} & \\
\textbf{Outcome} & \textbf{(4--8)} & \textbf{(8--15)} & \textbf{(15--32)} & \textbf{(32--82)} & \textbf{(82--189)} & \textbf{RR} & \textbf{$P$} \\
\midrule
MI (\%) & 3.11 & 3.27 & 3.36 & 3.63 & 4.51 & \textbf{1.45} & $1.4 \times 10^{-45}$ \\
Aortic stenosis (\%) & 1.59 & 1.62 & 1.67 & 1.84 & 2.21 & \textbf{1.39} & $2.2 \times 10^{-18}$ \\
Carotid disease (\%) & 0.54 & 0.51 & 0.57 & 0.68 & 0.80 & \textbf{1.48} & $9.3 \times 10^{-10}$ \\
PVD (\%) & 3.46 & 3.51 & 3.52 & 3.71 & 4.11 & \textbf{1.19} & $5.0 \times 10^{-11}$ \\
Stroke (\%) & 2.30 & 2.27 & 2.40 & 2.40 & 2.51 & 1.09 & $1.1 \times 10^{-2}$ \\
TIA (\%) & 1.20 & 1.15 & 1.20 & 1.17 & 1.28 & 1.07 & 0.16 \\
\midrule
Composite ASCVD (\%) & 11.37 & 11.26 & 11.62 & 12.47 & 13.94 & \textbf{1.23} & $4.9 \times 10^{-50}$ \\
\bottomrule
\end{tabular}
\end{table}

% --- TABLE 3: Correlation Matrix ---
\begin{table}[H]
\centering
\caption{\textbf{Lp(a) correlation matrix in FH and non-FH populations.} Spearman $\rho$ with significance: * $P < 0.05$, ** $P < 0.01$, *** $P < 0.001$. Lp(a) is orthogonal to all traditional cardiovascular risk factors in FH, demonstrating that it operates as an entirely independent risk axis.}
\label{tab:correlations}
\small
\begin{tabular}{@{}lcccc@{}}
\toprule
& \multicolumn{2}{c}{\textbf{FH ($n = 2{,}370$)}} & \multicolumn{2}{c}{\textbf{Non-FH ($n = 372{,}830$)}} \\
\cmidrule(lr){2-3} \cmidrule(lr){4-5}
\textbf{Variable} & $\rho$ & $P$ & $\rho$ & $P$ \\
\midrule
ApoB & 0.029 & 0.31 & 0.080*** & $<$0.001 \\
LDL-C (direct) & 0.031 & 0.27 & 0.081*** & $<$0.001 \\
LDL-C (Friedewald) & 0.045 & 0.12 & 0.095*** & $<$0.001 \\
Total cholesterol & 0.007 & 0.81 & 0.069*** & $<$0.001 \\
HDL-C & $-$0.006 & 0.83 & 0.021*** & $<$0.001 \\
Triglycerides & $-$0.083 & 0.07 & $-$0.048*** & $<$0.001 \\
ApoB/LDL ratio & 0.005 & 0.81 & $-$0.010* & 0.05 \\
BMI & $-$0.001 & 0.97 & --- & --- \\
CRP & $-$0.031 & 0.27 & --- & --- \\
\bottomrule
\end{tabular}
\end{table}

% --- TABLE 4: CMR by Lp(a) quintile and aortic stenosis ---
\begin{table}[H]
\centering
\caption{\textbf{Cardiac MRI parameters by Lp(a) quintile and aortic stenosis status (UK Biobank, $n = 60{,}947$).} LVEF = left ventricular ejection fraction; LVEDV = LV end-diastolic volume; GLS = global longitudinal strain. Aortic stenosis is associated with elevated Lp(a), increased LV mass (pressure overload), and reduced LVEF, while Lp(a) alone does not impair myocardial function.}
\label{tab:cmr}
\small
\begin{tabular}{@{}lccccc@{}}
\toprule
\textbf{Parameter} & \textbf{Q1} & \textbf{Q3} & \textbf{Q5} & \textbf{$\rho$} & \textbf{$P$} \\
\midrule
\multicolumn{6}{l}{\textit{Lp(a) quintile gradient ($n = 60{,}947$)}} \\
LVEF (\%) & 59.6 & 59.7 & 59.5 & 0.003 & 0.49 \\
LV mass (g) & 86.0 & 85.0 & 86.0 & $-$0.013 & 0.001 \\
LVEDV (mL) & 146.6 & 145.1 & 146.7 & --- & --- \\
GLS (--\%) & 5.7 & 5.7 & 5.7 & $-$0.004 & 0.35 \\
ASCVD (\%) & 6.6 & 7.2 & 9.2 & --- & $<$0.001 \\
\midrule
\multicolumn{6}{l}{\textit{Aortic stenosis ($n = 428$) versus no AS ($n = 60{,}519$)}} \\
& \textbf{AS} & \textbf{No AS} & & & \textbf{$P$} \\
Lp(a) (nmol/L) & 54.4 & 43.9 & & & 0.0005 \\
LV mass (g) & 107.1 & 85.3 & & & $<$0.0001 \\
LVEF (\%) & 58.2 & 59.6 & & & 0.006 \\
\bottomrule
\end{tabular}
\end{table}

% ============================================================================
% EXTENDED DATA
% ============================================================================
\clearpage
\section*{Extended Data}

\textbf{Extended Data Figure 1.} FoldX thermodynamic landscape of all 202 LDLR missense variants identified in the UK Biobank, coloured by domain.

\textbf{Extended Data Figure 2.} AlphaMissense pathogenicity scores for all 359 LDLR variants in the Structural Severity Score (SSS) catalogue, validated against ClinVar (AUC = 0.96) and Islam \textit{et al.} functional assays (AUC = 0.76).

\textbf{Extended Data Figure 3.} VUS reclassification: 655 LDLR variants of uncertain significance reclassified as likely pathogenic ($n = 255$, 38.9\%), likely benign ($n = 225$, 34.4\%), or remaining VUS ($n = 142$, 21.7\%) using the SSS framework.

\textbf{Extended Data Figure 4.} Lp(a)-cholesterol correction analysis: waterfall plot showing the shift in ApoB/LDL ratio from uncorrected (0.289) to Lp(a)-corrected (0.300) values.

\textbf{Extended Data Figure 5.} Genotype proxy for \textit{LPA} genetic architecture: Lp(a) concentration categories as proxy for \textit{LPA} genotype, showing identical distributions in FH versus non-FH carriers, confirming independence of \textit{LPA} and \textit{LDLR} genetic loci.

\textbf{Extended Data Table 1.} Baseline characteristics of all three study cohorts.

\textbf{Extended Data Table 2.} Wales FH Registry: Lp(a) threshold analysis with relative risks and 95\% CI.

\textbf{Extended Data Table 3.} Cardiac MRI parameters: FH carriers ($n = 369$) versus non-FH ($n = 60{,}578$).

% ============================================================================
% REFERENCES
% ============================================================================
\clearpage
\section*{References}

\begin{thebibliography}{40}

\bibitem{brown1986} Brown, M.S. \& Goldstein, J.L. A receptor-mediated pathway for cholesterol homeostasis. \textit{Science} \textbf{232}, 34--47 (1986).

\bibitem{ctt2010} Cholesterol Treatment Trialists' Collaboration. Efficacy and safety of more intensive lowering of LDL cholesterol: a meta-analysis of data from 170,000 participants in 26 randomised trials. \textit{Lancet} \textbf{376}, 1670--1681 (2010).

\bibitem{sabatine2017} Sabatine, M.S. \textit{et al.} Evolocumab and clinical outcomes in patients with cardiovascular disease. \textit{N. Engl. J. Med.} \textbf{376}, 1713--1722 (2017). [FOURIER]

\bibitem{berg1963} Berg, K. A new serum type system in man---the Lp system. \textit{Acta Pathol. Microbiol. Scand.} \textbf{59}, 369--382 (1963).

\bibitem{mclean1987} McLean, J.W. \textit{et al.} cDNA sequence of human apolipoprotein(a) is homologous to plasminogen. \textit{Nature} \textbf{330}, 132--137 (1987).

\bibitem{clarke2009} Clarke, R. \textit{et al.} Genetic variants associated with Lp(a) lipoprotein level and coronary disease. \textit{N. Engl. J. Med.} \textbf{361}, 2518--2528 (2009).

\bibitem{burgess2018} Burgess, S. \textit{et al.} Association of LPA variants with risk of coronary disease and the implications for lipoprotein(a)-lowering therapies: a Mendelian randomization analysis. \textit{JAMA Cardiol.} \textbf{3}, 619--627 (2018).

\bibitem{nordestgaard2010} Nordestgaard, B.G. \textit{et al.} Lipoprotein(a) as a cardiovascular risk factor: current status. \textit{Eur. Heart J.} \textbf{31}, 2844--2853 (2010).

\bibitem{arsenault2020} Arsenault, B.J. \textit{et al.} Lipoprotein(a) levels, genotype, and incident aortic valve stenosis: a prospective Mendelian randomization study and bias-adjusted meta-analysis. \textit{Circ. Genomic Precis. Med.} \textbf{13}, e002711 (2020).

\bibitem{kamstrup2009} Kamstrup, P.R., Tybj{\ae}rg-Hansen, A. \& Nordestgaard, B.G. Extreme lipoprotein(a) levels and improved cardiovascular risk prediction. \textit{J. Am. Coll. Cardiol.} \textbf{61}, 1146--1156 (2013).

\bibitem{wilson2022} Wilson, D.P. \textit{et al.} Use of lipoprotein(a) in clinical practice: a biomarker whose time has come. A scientific statement from the National Lipid Association. \textit{J. Clin. Lipidol.} \textbf{16}, e1--e34 (2022).

\bibitem{tsimikas2020} Tsimikas, S. A test in context: lipoprotein(a): diagnosis, prognosis, controversies, and emerging therapies. \textit{J. Am. Coll. Cardiol.} \textbf{69}, 692--711 (2017).

\bibitem{khan2025} Khan, M. \textit{et al.} Genome-wide CRISPR screen identifies LDLR as primary mediator of Lp(a) cellular uptake. \textit{Nature} (2025). [Preprint/in press]

\bibitem{tsimikas2020statin} Tsimikas, S. \textit{et al.} Statin therapy increases lipoprotein(a) levels. \textit{Eur. Heart J.} \textbf{41}, 2275--2284 (2020).

\bibitem{oconnor2015} O'Donoghue, M.L. \textit{et al.} Lipoprotein(a), PCSK9 inhibition, and cardiovascular risk. \textit{Circulation} \textbf{139}, 1483--1492 (2019).

\bibitem{raal2015} Raal, F.J. \textit{et al.} Inhibition of PCSK9 with evolocumab in homozygous familial hypercholesterolaemia (TESLA Part B): a randomised, double-blind, placebo-controlled trial. \textit{Lancet} \textbf{385}, 341--350 (2015).

\bibitem{sharma2017} Sharma, M. \textit{et al.} Inhibition of lipoprotein(a) uptake by plasminogen and LDL receptor-related protein. \textit{J. Lipid Res.} \textbf{58}, 1103--1113 (2017).

\bibitem{romagnuolo2015} Romagnuolo, R. \textit{et al.} Lipoprotein(a) catabolism is regulated by proprotein convertase subtilisin/kexin type 9 through the low density lipoprotein receptor. \textit{J. Biol. Chem.} \textbf{290}, 11649--11662 (2015).

\bibitem{bhatt2020} Bhatt, D.L. Lipoprotein(a): a new therapeutic frontier. \textit{JAMA} \textbf{324}, 234--235 (2020).

\bibitem{yang2019} Yang, X.-P. \textit{et al.} Scavenger receptor-BI-mediated uptake of lipoprotein(a). \textit{J. Lipid Res.} \textbf{60}, 1120--1128 (2019).

\bibitem{deo2022} Deo, R.C. \textit{et al.} Genetic ancestry-associated differences in lipoprotein(a) concentration and cardiovascular risk. \textit{JAMA} \textbf{327}, 1547--1556 (2022).

\bibitem{coassin2019} Coassin, S. \textit{et al.} A comprehensive map of single-base polymorphisms in the hypervariable LPA kringle IV type 2 copy number variation region. \textit{J. Lipid Res.} \textbf{60}, 1744--1756 (2019).

\bibitem{virani2012} Virani, S.S. \textit{et al.} Associations between lipoprotein(a) levels and cardiovascular outcomes in Black and White subjects. \textit{Circulation} \textbf{125}, 241--249 (2012).

\bibitem{patel2021} Patel, A.P. \textit{et al.} Lp(a) (Lipoprotein[a]) concentrations and incident atherosclerotic cardiovascular disease: new insights from a large national biobank. \textit{Arterioscler. Thromb. Vasc. Biol.} \textbf{41}, 465--474 (2021).

\bibitem{tsimikas2005} Tsimikas, S. \textit{et al.} Oxidized phospholipids, Lp(a) lipoprotein, and coronary artery disease. \textit{N. Engl. J. Med.} \textbf{353}, 46--57 (2005).

\bibitem{boffa2019} Boffa, M.B. \& Koschinsky, M.L. Oxidized phospholipids as a unifying theory for lipoprotein(a) and cardiovascular disease. \textit{Nat. Rev. Cardiol.} \textbf{16}, 305--318 (2019).

\bibitem{abramson2024} Abramson, J. \textit{et al.} Accurate structure prediction of biomolecular interactions with AlphaFold 3. \textit{Nature} \textbf{630}, 493--500 (2024).

\bibitem{kwon2008} Kwon, H.J. \textit{et al.} Molecular basis for LDL receptor recognition by PCSK9. \textit{Proc. Natl Acad. Sci. USA} \textbf{105}, 1820--1825 (2008).

\bibitem{miles2021} Miles, L.A. \textit{et al.} Plasminogen receptors. \textit{J. Biomed. Biotechnol.} \textbf{2021}, 1--16 (2021).

\bibitem{rudenko2002} Rudenko, G. \textit{et al.} Structure of the LDL receptor extracellular domain at endosomal pH. \textit{Science} \textbf{298}, 2353--2358 (2002).

\bibitem{jeon2001} Jeon, H. \textit{et al.} Implications for familial hypercholesterolemia from the structure of the LDL receptor YWTD-EGF domain pair. \textit{Nat. Struct. Biol.} \textbf{8}, 499--504 (2001).

\bibitem{genedy2026} Genedy, N. \textit{et al.} Structural Severity Score for LDLR variants in familial hypercholesterolaemia: from protein structure to clinical outcome. [Manuscript in preparation] (2026).

\bibitem{sijbrands2004} Sijbrands, E.J.G. \textit{et al.} Mortality over two centuries in large pedigree with familial hypercholesterolaemia: family tree mortality study. \textit{BMJ} \textbf{322}, 1019--1023 (2001).

\bibitem{thanassoulis2013} Thanassoulis, G. \textit{et al.} Genetic associations with valvular calcification and aortic stenosis. \textit{N. Engl. J. Med.} \textbf{368}, 503--512 (2013).

\bibitem{capoulade2020} Capoulade, R. \textit{et al.} Oxidized phospholipids, lipoprotein(a), and progression of calcific aortic valve stenosis. \textit{J. Am. Coll. Cardiol.} \textbf{76}, 1881--1890 (2020).

\bibitem{suk2006} Suk Danik, J. \textit{et al.} Influence of genetic and environmental determinants of lipoprotein(a) concentrations. \textit{Clin. Chem.} \textbf{52}, 138--146 (2006).

\bibitem{dahlen1990} Dahl{\'e}n, G.H. Lp(a) lipoprotein in cardiovascular disease. \textit{Atherosclerosis} \textbf{108}, 111--126 (1994).

\bibitem{sudlow2015} Sudlow, C. \textit{et al.} UK Biobank: an open access resource for identifying the causes of a wide range of complex diseases of middle and old age. \textit{PLoS Med.} \textbf{12}, e1001779 (2015).

\bibitem{tsimikas2020pelacarsen} Tsimikas, S. \textit{et al.} Lipoprotein(a) reduction in persons with cardiovascular disease. \textit{N. Engl. J. Med.} \textbf{382}, 244--255 (2020). [Pelacarsen Phase 2]

\bibitem{ohearn2023} O'Donoghue, M.L. \textit{et al.} Olpasiran for treatment of elevated lipoprotein(a): the OCEAN(a)-DOSE randomized clinical trial. \textit{JAMA} \textbf{330}, 725--735 (2023).

\bibitem{nissen2024} Nissen, S.E. \textit{et al.} Lepodisiran, an extended-duration short interfering RNA targeting lipoprotein(a): a randomized dose-ascending clinical trial. \textit{JAMA} \textbf{330}, 2075--2083 (2023).

\bibitem{nicholls2023} Nicholls, S.J. \textit{et al.} Muvalaplin, an oral small molecule inhibitor of lipoprotein(a) formation: a randomized clinical trial. \textit{JAMA} \textbf{330}, 1042--1053 (2023).

\end{thebibliography}

\end{document}
